\chapter{入出力とファイル}
\label{chap:io}

本章では、Python によるファイルの開閉、複数ファイルの探索、行単位の読み込み、値への変換を扱う。入力形式、文字コード、列型、時刻と単位、出力先を記録することで、同じデータ変換を再実行できる入出力\index{にゅうしゅつりょく@入出力|textbf}を構成する。

\section{テキストファイルの入出力と探索}

\subsection{ファイルのオープンと基本的な読み書き}

Python でテキストファイルを読むときの基本は \code{open}\index{open@\code{open}|textbf} である。さらに \code{with} を組み合わせると、処理が正常に終わった場合だけでなく、途中で例外が発生した場合にもファイルが閉じられる。例として、\code{sample.txt} に次の 3 行が入っているとする。

\begin{lstlisting}[numbers=none, xleftmargin=2\zw, title={sample.txt}]
run 1: 12.4
run 2: 10.1
run 3: 14.8
\end{lstlisting}

このファイルを 1 行ずつ表示するコードは次のようになる。

\begin{lstlisting}[language=Python]
path = "sample.txt"

with open(path, "r", encoding="utf-8") as f:
    for line in f:
        print(line.rstrip("\r\n"))
\end{lstlisting}

実行結果は次の通りである。

\begin{lstlisting}[numbers=none, xleftmargin=2\zw]
run 1: 12.4
run 2: 10.1
run 3: 14.8
\end{lstlisting}

第 1 引数の \code{path} は開くファイルの場所、\code{"r"} は読み込みモード、\code{encoding="utf-8"} は文字を UTF-8 として解釈する指定である。\code{as f} によって、開いたファイルを表すオブジェクトを変数 \code{f} で参照する。\code{for line in f} は全行を一度にメモリへ置かず、1 行ずつ取り出す。

引数なしの \code{read()} は、現在の読み取り位置から末尾までを一度に読み込む。開いた直後ならファイル全体が対象になる。入力容量に比例したメモリを要するため、大きなテキストファイルは反復処理によって 1 行ずつ読み込む。

このとき、\code{line} には行末文字が含まれる。\verb|rstrip("\r\n")| は CR と LF だけを末尾から取り除き、データに含まれる通常の空白は残す。引数なしの \code{rstrip()} は末尾の空白も削除するため、空白自体がデータである場合には区別が必要である。

\subsubsection{\code{pathlib} によるパス操作}

\code{pathlib.Path}\index{pathlib@\code{pathlib}|textbf} は、パス要素の結合、存在確認、名前や拡張子の取得をメソッドと演算子で提供する。

\begin{lstlisting}[language=Python]
from pathlib import Path

root = Path(".")
path = root / "sample.txt"

print(path.exists())
print(path.name)
\end{lstlisting}

この例でファイルが存在する場合、出力は次のようになる。

\begin{lstlisting}[numbers=none, xleftmargin=2\zw]
True
sample.txt
\end{lstlisting}

\code{exists()} はパスが存在するかを真偽値で返し、\code{name} は末尾のファイル名だけを返す。\code{root / "sample.txt"} の \code{/} は除算ではなく、\code{Path} オブジェクトに対してパス要素を結合する演算子である。例えば \code{Path("data") / "2026" / "04" / "sample.txt"} とすれば、複数階層のパスも同じ方法で組み立てられる。

\code{Path} の \code{/} 演算子は、OS ごとの区切り文字を直接記述せずにパス要素を結合する。\code{name}、\code{parent}、\code{with\_suffix()} は、それぞれ末尾名、親ディレクトリ、拡張子の置換を表す。\code{pathlib} を用いて出力ディレクトリとファイルパスを作る例を次に示す。

\begin{lstlisting}[language=Python]
outdir = Path("output")
outdir.mkdir(parents=True, exist_ok=True)
outfile = outdir / "summary.txt"
\end{lstlisting}

\code{parents=True} は \code{output} より上位の未作成ディレクトリも必要に応じて作る指定であり、\code{exist\_ok=True} は既存のディレクトリをエラーにしない指定である。ただし、同名の通常ファイルが存在する場合はエラーになる。

この指定により、出力ディレクトリが存在しない場合は作成し、既存のディレクトリであれば処理を継続する。出力ファイルの上書き可否は、\code{mkdir} とは別に検査する。

\subsubsection{ファイルオブジェクトのメソッド}

ファイルオブジェクトは、\code{read}、\code{readline}、\code{write} などのメソッドを持つ。ファイル全体を文字列として取得する処理と、現在位置から 1 行を取得する処理を区別する。Python 公式チュートリアル~\cite{python_tutorial} に基づく最小例を次に示す。

\begin{lstlisting}[language=Python]
with open(path, "r", encoding="utf-8") as f:
    first_line = f.readline()

with open("out.txt", "w", encoding="utf-8") as f:
    f.write("hello\n")
\end{lstlisting}

\code{readline()} の返り値には末尾の改行が含まれる。\code{write()} は書き込んだ文字数を返すが、改行を自動では補わないため、行末が必要なら \verb|"\n"| を明示する。この例を実行すると、\code{out.txt} の内容は \code{hello} と改行になる。

\code{readline} は現在位置から 1 行を読み、\code{write} は指定した文字列をファイルへ書く。大きな入力では、全体を読み込む前に \code{readline} や第~\ref{chap:shell}~章のシェルコマンドで先頭数行を確認できる。

\subsubsection{複数ファイルの収集}

ログや測定結果が複数ファイルに分かれている場合は、処理前に対象パスの一覧を確定する。\code{glob.glob}\index{glob@\code{glob}|textbf} による再帰的な探索例を次に示す。

\begin{lstlisting}[language=Python]
import glob

paths = sorted(glob.glob("data/**/*.txt", recursive=True))
print(f"count = {len(paths)}")
for path in paths[:3]:
    print(path)
\end{lstlisting}

例えば対象が 3 ファイルなら、出力は次のようになる。

\begin{lstlisting}[numbers=none, xleftmargin=2\zw]
count = 3
data/2025/run01.txt
data/2026/run01.txt
data/2026/run02.txt
\end{lstlisting}

パターン中の \code{*} は同じ階層の任意の文字列、\code{**} は \code{recursive=True} と組み合わせたときに複数階層へ一致する。\code{paths[:3]} は確認のために先頭 3 件だけを表示するスライスであり、処理対象を 3 件へ制限しているわけではない。

\code{glob.glob} はシェルのワイルドカードに近い書式で使える。複数ファイルを処理する前に、見つかった件数と先頭数件のパスを表示すれば、探索範囲の誤りを検出しやすい。

\code{glob.glob} の返す順序を処理条件として仮定せず、時系列や連番を順に処理する場合は \code{sorted} で順序を明示する。これにより、処理順序とログの順序を実行環境に依存させずに再現できる。

\subsection{Python での標準入出力}
\index{sys@\code{sys}|textbf}\index{ひょうじゅんにゅうりょく@標準入力!Python}\index{ひょうじゅんしゅつりょく@標準出力!Python}\index{ひょうじゅんえらーしゅつりょく@標準エラー出力!Python}
Python では、\code{sys.stdin}、\code{sys.stdout}、\code{sys.stderr} が三つの標準的な経路に対応する~\cite{python_sys}。\code{print()} は既定で標準出力へ書き込み、\code{file=} で別の経路を指定できる。次の内容を \code{uppercase.py} に保存する。
\begin{lstlisting}[language=Python]
import sys

n_lines = 0
for line in sys.stdin:
    print(line.rstrip("\r\n").upper())
    n_lines += 1
print(f"Processed {n_lines} lines", file=sys.stderr)
\end{lstlisting}
\begin{lstlisting}[language=bash]
python uppercase.py < sample.txt > converted.txt 2> conversion.log
\end{lstlisting}
変換した本文は \code{converted.txt} に、処理行数は \code{conversion.log} に保存される。\code{input()} は標準入力から 1 行を読み、末尾の改行を除いた文字列を返す関数である。引数に案内文を渡すと標準出力へ表示するため、データをパイプで受け渡すプログラムでは案内文の混入に注意する。

\section{文字列データの解析と整形}

\subsection{CSV 規則に従った列名付きデータの読み込み}\label{sec:csv-reader}

CSV\index{csv@CSV|textbf} では、引用符で囲まれたフィールド内にカンマや改行を含められる。\code{line.split(",")} はこの規則を解釈できないため、標準ライブラリの \code{csv} でレコードとフィールドを読み取る。

\begin{lstlisting}[language=Python]
import csv
from pathlib import Path

path = Path("events.csv")

with path.open("r", encoding="utf-8", newline="") as f:
    reader = csv.DictReader(f)
    first_row = next(reader)

print(first_row["time"])
print(first_row["mag"])
\end{lstlisting}

ここで \code{events.csv} の先頭が次の内容だとする。

\begin{lstlisting}[numbers=none, xleftmargin=2\zw, title={events.csv}]
time,mag,place
2026-04-04T12:00:00Z,3.2,"Offshore, Japan"
2026-04-04T12:05:00Z,2.8,Tokyo
\end{lstlisting}

先頭データ行の \code{time} と \code{mag} が文字列として読み込まれるため、実行結果は次の通りである。

\begin{lstlisting}[numbers=none, xleftmargin=2\zw]
2026-04-04T12:00:00Z
3.2
\end{lstlisting}

\code{DictReader}\index{dictreader@\code{DictReader}|textbf} は 1 行目を列名として読み、その後の各行を辞書として返す。したがって \code{row["mag"]} のように列名で値を取り出せる。\code{newline=""} は、テキスト層で改行変換を行わず、CSV モジュールに改行処理を任せる指定である。

列内のカンマが CSV の規則に従って引用されていれば、\code{csv} モジュールは一つのフィールドとして返す。上の \code{place} 列の \code{"Offshore, Japan"} も分割されない。

\subsection{文字列から値への変換と結果の出力}

多くの表形式データは、1 行の中に複数の列が並ぶ形で保存されている。テキストとして読んだだけでは、まだ単なる文字列である。そこから列に分け、整数や浮動小数点数へ変換して初めて解析に使える。たとえば空白区切りの 1 行を分解する例は次のようになる。

\begin{lstlisting}[language=Python]
line = "17 2026-04-04 12:34:56 15.2 OK"
columns = line.split()

record_id = int(columns[0])
date_text = columns[1]
time_text = columns[2]
value = float(columns[3])
status = columns[4]

print(record_id, date_text, time_text, value, status)
\end{lstlisting}

\code{split()} は引数を省略すると連続する空白を一つの区切りとして扱う。\code{int()} と \code{float()} が文字列を数値へ変換するため、出力は次のようになる。

\begin{lstlisting}[numbers=none, xleftmargin=2\zw]
17 2026-04-04 12:34:56 15.2 OK
\end{lstlisting}

入力仕様に基づいて、各列の型と単位を決める。識別子の先頭の 0 や、日時文字列の元表現を保持する必要がある場合は、数値へ変換せず文字列として扱う。変換処理を読み込み直後の一か所にまとめると、文字列と数値が混在する範囲を限定できる。

\subsubsection{ファイルへの書き出し}

複数行のテキストを出力ファイルへ順に書き込む例を示す。

\begin{lstlisting}[language=Python]
lines = ["run_id count mean", "1 120 3.42", "2 118 3.37"]

with open("summary.txt", "w", encoding="utf-8") as f:
    for line in lines:
        f.write(line + "\n")
\end{lstlisting}

このコードの実行後、\code{summary.txt} には \code{lines} の 3 要素が 1 行ずつ保存される。\code{"w"} は、ファイルが存在すれば内容を空にしてから書き、存在しなければ新しく作るモードである。これに対して \code{"a"} は既存内容の末尾へ追記する。既存ファイルを上書きしてよいのか、別名で保存すべきかは、実験ノートやレポートの再現性にも関わる。出力ファイル名の設計も解析手順に含まれる。

\subsubsection{文字列整形}

f-string（フォーマット済み文字列リテラル）は、文字列内に値と表示書式を埋め込む。

\begin{lstlisting}[language=Python]
mean = 3.14159
sigma = 0.27182

text = f"mean = {mean:.3f}, sigma = {sigma:.3f}"
print(text)
\end{lstlisting}

実行結果は \code{mean = 3.142, sigma = 0.272} となる。\code{f} は固定小数点形式、\code{.3} は小数点以下の桁数を指定する。これは表示上の丸めであり、変数 \code{mean} と \code{sigma} に保存された値そのものは変化しない。

\code{:.3f} は小数点以下 3 桁の固定小数点形式を指定する。表示桁数は、測定精度や不確かさに基づいて決め、表、図、本文で同じ丸め規則を用いる。

Python 公式チュートリアル~\cite{python_tutorial} では、f-string のほかに \code{str.format} や \code{\%} 演算子による書式指定も紹介されている。本書では、式と書式を同じ文字列内に明示できる f-string を主に用いる。

\section{オブジェクトと構造化データの保存}

\subsection{JSON と Pickle / PKL}

結果や設定を構造化して保存したいときは、JSON\index{json@JSON|textbf} を利用できる。テキストとして内容を確認でき、Python では標準ライブラリ \code{json} で扱える。保存と読み込みの基本形を次に示す。

\begin{lstlisting}[language=Python]
import json

config = {"threshold": 3.5, "bins": 40}

with open("config.json", "w", encoding="utf-8") as f:
    json.dump(config, f, indent=2)

with open("config.json", "r", encoding="utf-8") as f:
    loaded_config = json.load(f)

print(loaded_config)
\end{lstlisting}

\code{indent=2} は、入れ子の階層ごとに 2 個の空白で字下げし、要素ごとに改行した JSON を保存する指定である。読み戻した結果は次のように表示される。

\begin{lstlisting}[numbers=none, xleftmargin=2\zw]
{'threshold': 3.5, 'bins': 40}
\end{lstlisting}

JSON のオブジェクトは Python では辞書として復元される。

JSON は、キーと値、配列、入れ子構造をテキストとして保持する。一方、列ごとの型や圧縮方式を持たないため、大規模な表の反復読み込みには Parquet などの列指向形式も検討する。

\subsubsection{Pickle / PKL}

Python オブジェクトをそのまま保存する形式として、\code{pickle}\index{pickle@\code{pickle}|textbf} がある。拡張子には慣習的に \code{.pkl} が使われる。辞書、リスト、NumPy 配列、学習済みモデル、中間集計結果などを作業用に一時保存する例を次に示す。

\begin{lstlisting}[language=Python]
import pickle

result = {"mean": 3.14, "sigma": 0.27, "values": [1.0, 2.0, 3.0]}

with open("result.pkl", "wb") as f:
    pickle.dump(result, f)

with open("result.pkl", "rb") as f:
    loaded_result = pickle.load(f)

print(loaded_result["mean"])
print(loaded_result["values"])
\end{lstlisting}

復元後の出力は次の通りである。

\begin{lstlisting}[numbers=none, xleftmargin=2\zw]
3.14
[1.0, 2.0, 3.0]
\end{lstlisting}

\code{"wb"} と \code{"rb"} の \code{b} はバイナリモードを意味する。JSON と違って内容をテキストとして確認できないが、Python オブジェクトの構造を保存できる。ただし、別の Python やライブラリのバージョンで常に復元できる保証はなく、他言語との交換にも向かない。長期保存や共同研究の交換形式ではなく、生成環境と入力を記録したローカルな作業用キャッシュとして用いる。

Pickle の読み込みには、コード実行に関する注意も必要である。\code{pickle.load()} は、オブジェクトを復元する過程で任意のコードを実行する可能性がある。出所と生成過程を確認できない \code{.pkl} は読み込まず、自分で生成したもの、または信頼できる相手から安全な経路で受け取ったものだけを対象とする。

\section{日時データの変換と保存形式の選択}

\subsection{日付・時刻の数値表現}

日付や時刻を比較し、差を計算する場合は、文字列を時刻オブジェクトまたは数値表現へ変換する。標準ライブラリの \code{datetime}\index{datetime@\code{datetime}|textbf} を用いた例を次に示す。

\begin{lstlisting}[language=Python]
from datetime import datetime, timezone

text = "2026-04-04 12:34:56.123456"
dt = datetime.strptime(text, "%Y-%m-%d %H:%M:%S.%f")
# 入力の時刻が UTC であることを確認した上で指定する
dt = dt.replace(tzinfo=timezone.utc)
unix_time = dt.timestamp()
print(dt)
print(unix_time)
\end{lstlisting}

\code{strptime()} の第 1 引数は変換する文字列、第 2 引数は文字列の形式である。\code{\%Y}、\code{\%m}、\code{\%d} は年、月、日を表し、\code{\%H}、\code{\%M}、\code{\%S} は時、分、秒、\code{\%f} はマイクロ秒を表す。\code{timestamp()} は Unix epoch（1970 年 1 月 1 日 00:00:00 UTC）からの秒数を返すが、タイムゾーン情報を持たない \code{datetime} の解釈は実行環境のローカルタイムゾーンに依存する。UTC のログでは、タイムゾーンを明示するか、後述する Astropy の \code{Time} を使う必要がある。

時刻差を保持する単位は、必要な時間分解能と値域に基づいて選ぶ。別の表現へ変換する場合は、基準時刻、時刻系、単位、丸め方法を記録する。

\code{datetime} オブジェクトは暦に沿った表示と計算を担う。一方、大量の時刻差を配列演算で求める場合は、基準時刻からの秒やナノ秒として表す方が適することもある。表示用の表現と計算用の表現を分ける場合は、基準時刻、単位、時刻系を記録する。

さらに、タイムゾーンが関わるデータでは、見かけの時刻だけを比較すると誤解が生じる。研究用のログや観測データでは UTC を使うことが多いが、比較前に各データの基準時刻を確認する必要がある。


\subsection{CSV・JSON とバイナリ形式の選択}

CSV はテキストとして内容を確認できる一方、読み込み時には文字列から各列型への変換が必要となる。列数が多い場合は、使用しない列の読み込みも所要時間とメモリ使用量を増やす。保存形式は、交換先が利用できる実装、読み込む列と行の範囲、測定した入出力時間を基準に選ぶ。

\subsubsection{CSV・JSON の可読性と変換コスト}
テキスト形式は、UTF-8 などの文字コードを用いて人間が読める形で保存する。CSV/TSV は多くのプログラミング言語や表計算ソフトで読み書きできるため、データ交換に利用しやすい。

計算に使う際には、ディスク上の文字列（例: 「1.2345」）を実数値へ変換する処理（パース）が必要になる。大規模データでは、この変換と区切り文字の解析が読み込み時間の無視できない部分を占めることがある。

\subsubsection{NPY・HDF5・Parquet の用途}
バイナリ形式は、数値を文字列から変換する処理を省ける場合があり、CSV より短時間で読み書きできることが多い。ただし、実際の時間はファイルサイズ、圧縮方式、記憶装置、列数、読み込む範囲、ライブラリの実装に左右される。バイナリ形式に変えただけで、常に記憶装置の転送速度だけが上限になるわけではない。

代表的なバイナリ形式には用途の違いがあり、データ構造と交換相手に合わせて選ぶ必要がある。

\begin{itemize}
    \item \textbf{NPY / NPZ (NumPy 標準)}: NumPy 配列の形と型を保持する。NPY は 1 配列、NPZ は複数の NPY データを ZIP コンテナへまとめる。交換先が NumPy 形式に対応しているかを確認する。
    \item \textbf{Pickle / PKL}: Python オブジェクトを直列化できるため、同じ管理環境内の一時的な中間結果やキャッシュに用いられる。Python 以外との交換には適さず、読み込み時に任意のコードが実行され得るため、信頼できないファイルを読み込んではならない。
    \item \textbf{HDF5}: 大規模な数値データを階層構造で管理できる形式である。多次元配列を部分的に読み書きし、データとメタデータを一つのファイルへまとめたい場合に向く。C/C++、Fortran、R などからも利用できるが、受渡し相手が使うライブラリとデータ構造を事前に確認する必要がある。
    \item \textbf{Parquet}: 列指向のバイナリ形式である。pandas や Polars から表を保存し、列選択や行グループ単位の読み込みを行える。レコード単位の更新ではなく、列の抽出、集計、複数ファイルの走査を前提とする。
\end{itemize}

解析の初期段階では、内容をテキストとして確認でき、複数のソフトウェアで扱える CSV が検証に役立つ。大量のイベントデータの反復読み込みが測定上のボトルネックになった場合は、変換手順と列型を記録したうえで Parquet や HDF5 を中間形式として導入する。小規模なデータや一度しか読み込まないデータでは、変換時間と形式管理の負担が利点を上回る場合がある。\code{.pkl} は、信頼できる同一環境内の途中結果や作業用キャッシュに用途を限定する。

Parquet は保存形式の名前であり、\code{pyarrow} は pandas が Parquet の読み書きに利用できる実装の一つである。\code{read\_parquet()} と \code{to\_parquet()} の実行には、pandas が対応するエンジンを利用できる環境が必要となる。

\section{章末課題：USGS カタログの読み込みと要約}

第~\ref{chap:shell}~章で USGS の CSV を取得していない場合は、\code{raw} ディレクトリを作り、次の内容を \code{raw/usgs\_example.csv} として保存すれば、同じ課題を小規模な入力で実行できる。その場合は課題中の入力パスを読み替える。人工データは、取得済みの実カタログとは別のファイルへ保存する。

\begin{lstlisting}[numbers=none, xleftmargin=2\zw, title={raw/usgs\_example.csv}]
time,latitude,longitude,mag,place
2024-01-01T00:00:00.000Z,35.0,140.0,3.2,"Offshore, Japan"
2024-01-01T00:05:00.000Z,36.0,141.0,,Japan
2024-01-01T00:12:30.000Z,34.5,139.5,2.8,Tokyo
\end{lstlisting}

この小規模入力では、ヘッダを除くデータ行数は 3、最初と最後の時刻はそれぞれ \code{2024-01-01T00:00:00.000Z} と \code{2024-01-01T00:12:30.000Z}、空欄を除くマグニチュードの最小値と最大値は 2.8 と 3.2 になる。自分の出力がこれらと一致することを確認した後、実データへ進め。

\begin{enumerate}
\item \textbf{USGS CSV の要約と JSON への保存}。
第~\ref{chap:shell}~章で取得した \code{raw/usgs\_global\_2024\_01.csv} を \code{csv.DictReader} で読み、少なくとも次の情報を \code{results/usgs\_summary.json} へ保存せよ。
\begin{itemize}
\item ヘッダを除くデータ行数
\item 最初のイベント時刻
\item 最後のイベント時刻
\item マグニチュードの最小値と最大値
\end{itemize}
空欄の \code{mag} は数値集計から除き、その件数も記録せよ。有効な値が一つもない場合は、最小値・最大値を \code{None} として保存する。出力には \code{pathlib} と \code{json.dump} を用いよ。
\code{Path.mkdir(parents=True, exist\_ok=True)} で出力ディレクトリを作り、\code{json.dump(..., indent=2)} の各引数が何を指定するかも説明せよ。保存後に \code{json.load()} で読み戻し、辞書の内容を表示して検証せよ。

\item \textbf{解析用中間データの Pickle への保存}。
CSV の先頭 200 行から \code{time}、\code{latitude}、\code{longitude}、\code{mag} を抽出し、辞書のリストとして \code{results/usgs\_preview.pkl} へ保存せよ。読み直して要素数と先頭 1 件を表示し、保存前の値と一致することを確かめよ。
小規模入力を使った場合、要素数は 3 であり、先頭辞書の \code{time} は \code{2024-01-01T00:00:00.000Z} になる。\code{"wb"} と \code{"rb"} の \code{b} の意味を説明し、この PKL を信頼できない相手から受け取った場合に読み込んではならない理由も述べよ。

\item \textbf{発展課題}。
CSV の中には空欄や想定外の値が入り得る。空欄に加えて数値へ変換できない \code{mag} もスキップし、何件スキップしたかを \code{results/skipped\_rows.txt} へ書き出せ。\code{DictReader} では空欄セルは空文字列として読まれるので、たとえば \code{if row["mag"] == ""} のような判定を入れればよい。
小規模入力に対するスキップ件数は 1 件となる。\code{float(row["mag"])} が失敗する場合も \code{ValueError} として数えられるようにし、空欄と不正な文字列を区別して記録せよ。
\end{enumerate}

ファイルの構造を確認し、必要な値を Python のオブジェクトへ変換できれば、多数の数値を同じ型の配列として計算できる。次章では、ここで読み込んだ時刻や信号値を NumPy 配列へ変換し、差分、条件抽出、集計、ヒストグラムを扱う。
